\documentclass[11pt]{article}
\usepackage{fontspec}
\usepackage[margin=1in,top=1.1in,bottom=1.1in]{geometry}
\usepackage{booktabs}
\usepackage{longtable}
\usepackage{array}
\usepackage{amsmath}
\usepackage{xcolor}
\usepackage{titlesec}
\usepackage{fancyhdr}
\usepackage{enumitem}
\usepackage{caption}
\usepackage{unicode-math}
\usepackage[hidelinks]{hyperref}
\usepackage{microtype}
\usepackage{tabularx}
\usepackage{colortbl}
\usepackage{needspace}
\usepackage{float}

\setmainfont{Times New Roman}
\setmonofont{Consolas}[Scale=0.86]
\setmathfont{Cambria Math}

\definecolor{dqblue}{HTML}{1F3864}
\definecolor{dqgrey}{HTML}{5A5A5A}
\definecolor{dqrule}{HTML}{B4B4B4}
\definecolor{rowtint}{HTML}{F2F2F2}
\definecolor{warnbg}{HTML}{FBF3E4}

\titleformat{\section}{\Large\bfseries\color{dqblue}}{\thesection}{0.7em}{}
\titleformat{\subsection}{\large\bfseries\color{dqblue}}{\thesubsection}{0.6em}{}
\titleformat{\subsubsection}{\normalsize\bfseries}{\thesubsubsection}{0.5em}{}
\titlespacing{\section}{0pt}{1.6ex plus .3ex}{0.8ex}
\titlespacing{\subsection}{0pt}{1.3ex plus .3ex}{0.6ex}

\setlength{\headheight}{21.3pt}
\addtolength{\topmargin}{-7.3pt}
\setlength{\emergencystretch}{3em}
\pagestyle{fancy}
\fancyhf{}
\fancyhead[L]{\footnotesize\color{dqgrey}DynQuant --- Signal-Driven Mixed-Precision Quantization}
\fancyhead[R]{\footnotesize\color{dqgrey}\thepage}
\renewcommand{\headrulewidth}{0.4pt}
\renewcommand{\headrule}{\hbox to\headwidth{\color{dqrule}\leaders\hrule height \headrulewidth\hfill}}

\captionsetup{font=small,labelfont={bf,color=dqblue}}
\setlist{itemsep=0.25ex,topsep=0.5ex,parsep=0pt}

\newcommand{\x}{$\times$}
\newcommand{\pt}{\,pts}
\newcommand{\pct}{\,\%}
\newcommand{\code}[1]{\texttt{#1}}
\newcommand{\minus}{$-$}

% keeps numeric columns aligned on the decimal point without dcolumn
\newcolumntype{R}{>{\raggedleft\arraybackslash}p{0pt}}

\begin{document}

%% ============================================================ TITLE
\begin{titlepage}
\centering
\vspace*{1.2in}
{\color{dqblue}\rule{\textwidth}{1.2pt}}\\[0.8em]
{\Huge\bfseries DynQuant\par}
\vspace{0.5em}
{\LARGE Signal-Driven Mixed-Precision Quantization\\for Large Language Models\par}
\vspace{0.7em}
{\color{dqblue}\rule{\textwidth}{1.2pt}}\\[2.0em]
{\large A Complete Technical White Paper\par}
\vspace{0.4em}
{\large Method, On-Disk Format, CUDA Runtime, and Measured Results\par}
\vspace{2.5em}

\begin{minipage}{0.86\textwidth}
\small
\textbf{Summary.} DynQuant assigns a different numeric width to every weight tensor in a
language model, chosen from what the model itself reveals about that tensor while it is being
fine-tuned. Two per-channel second moments harvested by a training callback --- the mean square
of each input activation channel and of each output-gradient channel --- combine with the
\emph{actual} quantization error of the \emph{actual} candidate encoding to give a per-module,
per-width sensitivity in units of loss. A greedy multiple-choice knapsack then buys bits where
they are cheapest per unit of loss avoided, subject to soft quality floors and a small set of
hard structural floors that no budget may breach. Weights are encoded group-wise with a float
offset and an MSE-optimal per-group clip search, packed into \code{int32} words whose boundaries
align to the group boundaries, and executed by CUDA kernels that never materialise the
full-precision tensor.
\\[0.7em]
On Qwen3.5-2B-Base fine-tuned for legal holding selection (CaseHOLD, $n=5314$), the technique
retains \textbf{92\pct} of the fine-tuning gain at an average of \textbf{3.25 bits} --- against
69\pct{} for the ranking score published in the original work and 73.79\pct{} absolute for a
uniform 3-bit control that it beats by \textbf{$+10.29$\pt} ($p = 1.3\times10^{-81}$, exact
paired McNemar). Peak allocated VRAM falls \textbf{4.84\x}, the packed model's accuracy is
\emph{bit-identical} to the simulated one, and 417 kernel-parity tests pass under all four
\code{compute-sanitizer} tools with zero errors and zero hazards.
\end{minipage}

\vfill
{\small\color{dqgrey}
Measured on NVIDIA A100 80\,GB PCIe (108 SMs, 40\,MB L2), CUDA 13.0.88,
PyTorch 2.13.0+cu130, Transformers 5.14.1.\\
Every number in this document is measured, not projected. Section~\ref{sec:notestablished}
states what is \emph{not} established.\par}
\vspace{0.6em}
\end{titlepage}

\tableofcontents
\newpage

%% ============================================================ 1
\section{Executive summary}

\subsection{What the technique does}

A quantized language model is a budget-allocation problem. You have a fixed number of bits to
spend across a few hundred weight tensors, and the tensors are not equally deserving: some
absorb aggressive rounding with no measurable effect on the loss, others are destroyed by it.
Uniform quantization ignores the difference. DynQuant measures it.

The measurement happens during fine-tuning, which is when the model is already doing the forward
and backward passes that expose the relevant quantities. A callback registered on the
\code{Trainer} accumulates, per linear module, two vectors of per-channel second moments: the
mean square of each input activation channel, $\mathbb{E}[x_c^2]$, of length \code{in\_features};
and the mean square of each output-gradient channel, $\mathbb{E}[\delta_r^2]$, of length
\code{out\_features}. Nothing else is needed. From those two vectors and the weight tensor
itself, the sensitivity of the loss to quantizing module $m$ at width $b$ is
\begin{equation}
\mathcal{S}_m(b) \;=\; \sum_{r,c}\ \mathbb{E}[\delta_r^2]\;\mathbb{E}[x_c^2]\;
\bigl(W_{rc} - Q_b(W)_{rc}\bigr)^2 ,
\label{eq:sens}
\end{equation}
a diagonal Gauss--Newton (equivalently, empirical-Fisher) estimate of the loss increase, in loss
units, computed by running the weight block through the \emph{real} encoder at the \emph{real}
candidate width, clip search included. There is no surrogate error curve and no rank
transformation. Section~\ref{sec:sens} derives it.

Allocation then becomes a multiple-choice knapsack over the widths $\{2,3,4,8\}$ where the price
of a move is the difference between two sensitivities --- what that bit actually buys:
$\mathcal{S}_m(3) - \mathcal{S}_m(4)$. Every module starts at its floor; if the floors overrun
the budget the allocator downgrades by cheapest damage; then it always runs an upgrade pass to
reclaim the slack that discrete width steps leave on the table. Section~\ref{sec:alloc}.

Encoding is group-wise asymmetric affine, group 128, with an unconstrained \emph{float} offset
rather than an integer zero-point, and a per-group MSE-optimal clip search whose objective is the
error of the real encoder--decoder round trip rather than of an idealised one.
Section~\ref{sec:encoder}.

Execution keeps the weights packed. \code{DynQuantLinear} holds \code{int32} buffers in VRAM and
dispatches to a CUDA GEMV templated on \code{<BITS, GROUP\_SIZE>} for $M \le 8$ rows, falling back
to tile dequant plus \code{F.linear} above that. Section~\ref{sec:runtime}.

\subsection{What changed from the published formulation}

The original work scored modules with an ordinal product of percentile ranks,
\begin{equation}
S_i \;=\; \operatorname{Rank}\!\bigl(\log(1 + \operatorname{Var}_t\lVert\nabla W_i\rVert^2)\bigr)
\;\times\;
\operatorname{Rank}\!\bigl(\operatorname{EMA}\lVert X_i\rVert\bigr),
\label{eq:oldscore}
\end{equation}
and priced a move as $\text{score} \times \text{params} \times \Delta\text{error}$ with
$\Delta\text{error}$ read off a universal $4^{-b}$ curve. Three things are wrong with that in
practice, and all three were confirmed by measurement on this model:

\begin{enumerate}
\item \textbf{The activation factor points the wrong way.} $\operatorname{EMA}\lVert X_i\rVert$
is a \emph{saliency} signal --- large activations mean the tensor is busy. Correlated against the
measured loss damage of quantizing each module, it comes out at $\rho = -0.2059$ globally and is
positive in only 2 of 12 roles. Large-activation modules on this architecture are the
\emph{robust} ones. Multiplying a useful signal by it is worse than dropping it.
\item \textbf{The rank product destroys magnitude.} Ordinal ranks answer ``which module is
busier'' but a knapsack needs ``by how much''. The shipped rank product correlates with measured
damage at $\rho = +0.0040$ --- statistically indistinguishable from noise.
\item \textbf{A universal error curve is not the error.} $4^{-b}$ is right on average and wrong
per tensor. Equation~\ref{eq:sens} substitutes the measured error of the encoding that will
actually be written to disk.
\end{enumerate}

Equation~\ref{eq:oldscore} is retained, unchanged, behind the \code{paper-3.15} preset, so
published numbers remain reproducible bit-for-bit. Equation~\ref{eq:sens} is the default.

\subsection{Headline results}

\begin{table}[H]
\centering\small
\caption{CaseHOLD exact match, Qwen3.5-2B-Base, $n=5314$. Chance is 20\pct{} (five-way holding
selection). ``Retention'' is the share of the $+53.35$\pt{} fine-tuning gain that survives
quantization.}
\begin{tabular}{lrrrr}
\toprule
Arm & Stored bits & Size (GiB) & Exact match & Retention \\
\midrule
Base model, fp16                       & 16.000 & 3.505 & 35.13\pct & --- \\
\rowcolor{rowtint}
Fine-tuned, fp16                       & 16.000 & 3.505 & \textbf{88.48\pct} & 100\pct \\
DynQuant $\approx$4.25 bits            & 4.249 & 0.931 & 86.49\pct & 96\pct \\
Uniform 4-bit control                  & 4.255 & 0.932 & 86.62\pct & 96\pct \\
\rowcolor{rowtint}
\textbf{DynQuant $\approx$3.25 bits, sensitivity} & 3.249 & 0.712 & \textbf{84.08\pct} & \textbf{92\pct} \\
Uniform 3-bit control                  & 3.256 & 0.713 & 73.79\pct & 73\pct \\
DynQuant $\approx$3.25 bits, rank product (Eq.~\ref{eq:oldscore}) & 3.249 & 0.712 & 71.75\pct & 69\pct \\
\midrule
Guessing                               & --- & --- & 20.00\pct & --- \\
\bottomrule
\end{tabular}
\end{table}

Three facts carry the paper. First, at 3.25 bits the sensitivity-driven allocation beats a
uniform 3-bit baseline by $+10.29$\pt{} and the published ranking score by $+12.33$\pt, both at
$p < 10^{-80}$ on the same 5314 problems. Second, the published ranking score is not merely
unhelpful at this budget --- it loses to uniform by $-2.03$\pt{} ($p = 5.8\times10^{-5}$), which
means the signal was actively misdirecting the budget. Third, at 4.25 bits neither allocation
separates from uniform ($-0.13$\pt, $p = 0.56$): the headroom to allocate only exists once the
budget is tight enough that some module must be hurt.

On the runtime side: peak allocated VRAM $0.7237$ GiB against $3.5052$ GiB in bf16, a
\textbf{4.84\x} reduction, with the allocator's predicted footprint landing within
\textbf{0.03\pct} of the live measurement; accuracy through the CUDA kernels
\emph{exactly} equal to the simulated path at $4468/5314$; worst observed relative error
$6.96\times10^{-3}$; and 417 parity tests green under memcheck, initcheck, racecheck and
synccheck.

And, stated plainly because it matters: decode throughput at batch 1 is
\textbf{0.90\x} of bf16, not faster. The kernel reaches 36--98\pct{} of achievable HBM bandwidth
depending on width, missing the $\ge$70\pct{} gate at 2, 3 and 4 bits.
Section~\ref{sec:notestablished} explains why, and why the reason is arithmetic rather than a
defect.

%% ============================================================ 2
\newpage
\section{The pipeline}

\begin{verbatim}
   +---------------------------------------------------------------------+
   | 1. FINE-TUNE          DynQuantCallback on transformers.Trainer      |
   |                       forward hook  -> E[x_c^2]   (in_features)     |
   |                       backward hook -> E[d_r^2]   (out_features)    |
   |                       sampled every channel_moment_every=16 steps   |
   |                       -> channel_moments.safetensors                |
   +---------------------------------------------------------------------+
                                     |
   +---------------------------------------------------------------------+
   | 2. CLASSIFY           graph/classify.py: structural inference over  |
   |                       the module tree -> ModuleRole per module,     |
   |                       RowPartition list for fused projections       |
   +---------------------------------------------------------------------+
                                     |
   +---------------------------------------------------------------------+
   | 3. SCORE              for each module m, for each b in {2,3,4,8}:   |
   |                       quantize -> measure error -> weight by the    |
   |                       two moment vectors -> S_m(b)  [loss units]    |
   +---------------------------------------------------------------------+
                                     |
   +---------------------------------------------------------------------+
   | 4. ALLOCATE           start at floors; greedy knapsack on           |
   |                       price = S_m(lo) - S_m(hi); downgrade if over  |
   |                       budget; ALWAYS upgrade to reclaim slack       |
   |                       -> bits per module, honest byte accounting    |
   +---------------------------------------------------------------------+
                                     |
   +---------------------------------------------------------------------+
   | 5. ENCODE             group 128 along in_features; float offset;    |
   |                       per-group clip search over 8 alphas scored    |
   |                       through the real codec; pack to int32 words   |
   |                       -> safetensors + manifest (versioned)         |
   +---------------------------------------------------------------------+
                                     |
   +---------------------------------------------------------------------+
   | 6. RUN                DynQuantLinear holds packed int32 in VRAM.    |
   |                       M <= 8   : gemv_nbit<BITS, GROUP_SIZE>        |
   |                       M >  8   : dequant tile -> F.linear           |
   +---------------------------------------------------------------------+
\end{verbatim}

Stages 1 and 6 touch the user's training and inference loops; stages 2--5 are offline and take
minutes. Nothing in stages 2--5 needs a GPU except for speed.

%% ============================================================ 3
\needspace{7\baselineskip}
\section{Signals collected during fine-tuning}
\label{sec:signals}

\subsection{What is collected, and why those quantities}

For a linear layer $y = Wx$ with $W \in \mathbb{R}^{\,\text{out}\times\text{in}}$, the gradient
of the loss with respect to the weight is the outer product $\nabla_W \mathcal{L} = \delta x^{\!\top}$
where $\delta = \partial\mathcal{L}/\partial y$. That identity is the whole reason the two moment
vectors suffice. A perturbation $\Delta W$ changes the loss, to second order, by
\begin{equation}
\Delta\mathcal{L} \;\approx\; \tfrac12 \operatorname{vec}(\Delta W)^{\!\top} H
\operatorname{vec}(\Delta W),
\end{equation}
and the Gauss--Newton approximation to $H$ for this layer factorises as
$H \approx \mathbb{E}[\delta\delta^{\!\top}] \otimes \mathbb{E}[xx^{\!\top}]$. Keeping only the
diagonal of each factor --- which is what makes the estimate affordable --- collapses the
quadratic form to exactly Equation~\ref{eq:sens}: an outer product of the two second-moment
vectors, contracted against the squared quantization error.

So the callback accumulates:

\begin{table}[H]
\centering\small
\caption{Signals accumulated per module by \code{SignalTracker}.}
\begin{tabular}{p{0.20\textwidth}p{0.13\textwidth}p{0.55\textwidth}}
\toprule
Quantity & Length & Accumulation \\
\midrule
$\mathbb{E}[x_c^2]$ & \code{in\_features} & Forward hook. Running mean of the squared input over
all tokens seen on sampled steps. Stored as \code{input\_sq}. \\
\rowcolor{rowtint}
$\mathbb{E}[\delta_r^2]$ & \code{out\_features} & Full backward hook on the output gradient.
Running mean of the square. Stored as \code{output\_grad\_sq}. \\
\code{observations} & scalar & Token count behind each mean, carried in safetensors metadata so
partial runs merge correctly. \\
\rowcolor{rowtint}
activation RMS EMA & scalar & Legacy saliency signal, $\beta = 0.99$. Retained for
\code{paper-3.15} and for diagnostics. \\
gradient variance & scalar & Legacy plasticity signal, Welford online variance of
$\lVert\nabla W\rVert^2$ across optimizer steps. \\
\bottomrule
\end{tabular}
\end{table}

\subsection{Cost, and the sampling gate}

The two reductions are cheap in FLOPs but they are extra kernel launches on tensors that are
already in cache, and there are hundreds of modules. Measured, collecting them on \emph{every}
step adds \textbf{2.30\pct} of step time --- which fits inside the 3\pct{} overhead budget, but
leaves nothing for anything else. Since the moments are statistical quantities whose estimates
change slowly, they are sampled: the \code{channel\_moment\_every} option defaults to 16, so the
hooks fire on one step in sixteen and the overhead falls by roughly that factor. The legacy
scalar signals are cheap enough to collect unconditionally.

All accumulators are device-resident tensors indexed by module id; there is no
\code{.cpu().item()} per module per step, which in the original research code caused roughly one
GPU synchronisation per module per step.

\subsection{Where the moments come from matters}

The moments are collected on data the model has \emph{not} memorised. Three splits were
available for CaseHOLD and only one is correct:

\begin{itemize}
\item \code{test} --- leaks the evaluation set into the allocation decision. Rejected.
\item \code{train} --- after fine-tuning, training loss is $0.0000$ to four decimals, so
$\mathbb{E}[\delta^2]$ collapses to $\sim 10^{-9}$ and the gradient factor carries no
information. Rejected on measurement, not on principle.
\item \code{validation} --- used. 512 rows, 204\,693 tokens. Split disjointness was verified:
$|\text{validation} \cap \text{test}| = 1$ out of 5314, a duplicate present in the source
dataset.
\end{itemize}

That the \code{train} split fails is worth dwelling on. A converged model has no gradient signal
on its training data, so any Fisher-style importance measure harvested there degenerates. The
signal must come from data the model still finds slightly surprising.

%% ============================================================ 4
\section{Cardinal sensitivity}
\label{sec:sens}

\subsection{The estimator}

\code{score/sensitivity.py} computes, for each module and each candidate width, the scalar of
Equation~\ref{eq:sens}. The implementation never forms the $[\text{out},\text{in}]$ weighted
product, which for the embedding table would be a 2\,GB intermediate. Instead it exploits that
the weighting factorises:

\begin{equation}
\mathcal{S}_m(b) \;=\; \underbrace{\Bigl[\,E^{\circ 2} \, x^2 \Bigr]}_{\text{length out}}
\cdot\; \delta^2 ,
\qquad E = W - Q_b(W),
\end{equation}
so a matrix--vector product against $\mathbb{E}[x^2]$ reduces the squared error to one number per
output row, and a dot product against $\mathbb{E}[\delta^2]$ finishes it. Rows are independent
because groups run along the input axis, so the computation is chunked by rows under a 64\,MiB
budget:

\begin{verbatim}
_ROW_CHUNK_BYTES = 64 << 20   # rows are independent (groups run along input axis)

def _module_sensitivity(weight, x2, d2, *, bits, group_size, symmetric) -> float:
    rows, cols = weight.shape
    per_row = max(cols * 4, 1)
    chunk = max(1, min(rows, _ROW_CHUNK_BYTES // per_row))
    total = 0.0
    for start in range(0, rows, chunk):
        stop  = min(start + chunk, rows)
        block = weight[start:stop]
        quantized, _ = quantize_with_search(block, bits=bits, group_size=group_size,
                                           symmetric=symmetric, row_offset=start)
        error = block.float() - quantized.dequantize(dtype=torch.float32)
        error.mul_(error)
        per_output = error @ x2f          # [chunk, in] @ [in] -> [chunk]
        total += float(per_output @ d2f[start:stop])
    return total
\end{verbatim}

The \code{row\_offset} argument matters: it makes the clip search deterministic and identical to
what the encoder will do on the full tensor, so the sensitivity is measured on the bytes that
will actually be written.

\subsection{What the estimate is, and what it is not}

\code{SensitivityTable} carries an \code{absolute: bool} flag, and it is \code{False} today.
The reason is honesty about calibration. The estimate is first-order in the module: it prices the
damage from quantizing module $m$ \emph{alone}, with everything else at full precision.
Summed over a whole allocation it overshoots the measured loss increase by roughly 2\x, because
errors in different modules partially cancel and because the quadratic model degrades at
2-bit-sized perturbations. So the table is used as \emph{a calibrated ordering with meaningful
proportions} --- module $A$ being three times more sensitive than module $B$ is a claim the
numbers support --- but not as a loss prediction. Nothing in the allocator needs the absolute
scale; the knapsack only compares prices.

Two implementation details that are easy to get wrong and are handled explicitly:

\begin{itemize}
\item \textbf{Missing moments must not read as zero.} A module with no usable moment vectors
(never routed, hook missed, shape disagreement) lands in \code{unestimable} and falls back to a
proxy score. Treating it as zero sensitivity would make it the first thing the allocator
destroys.
\item \textbf{Tied embeddings need aliasing.} On a tied model, \code{embed\_tokens.weight} and
\code{lm\_head.weight} are the same storage, but in the stored orientation only the head produces
a moment pairing. \code{\_tie\_aliases} resolves the pair. Separately, \code{\_shapes\_agree}
rejects wrong-axis broadcasts outright rather than letting NumPy-style broadcasting silently
produce a plausible-looking wrong number.
\end{itemize}

\subsection{The calibration-free control}

\code{weight\_only\_sensitivity} exists as the honest control: the same formula with both moment
vectors replaced by ones, i.e.\ pure weight reconstruction error. It requires no fine-tuning
data at all. It correlates with measured damage at $\rho = +0.0645$ globally, and in the policy
shootout of Section~\ref{sec:shootout} it recovers $-29.46\pct$ of the recoverable damage at
3.125 bits --- that is, it is \emph{worse than doing nothing}. This is the cleanest available
evidence that the activation and gradient statistics are load-bearing rather than decorative.

%% ============================================================ 5
\section{Allocation}
\label{sec:alloc}

\subsection{Move pricing}

The allocator's central object is the price of moving one module from width \code{lower} to
width \code{upper}:

\begin{verbatim}
def move_value(self, lower: int, upper: int) -> float:
    if self.sens is not None:
        lo, hi = self.sens.get(lower), self.sens.get(upper)
        if lo is not None and hi is not None:
            return max(lo - hi, 0.0)
    proxy = self.score * self.num_params * (_error_scale(lower) - _error_scale(upper))
    return max(self.fallback_scale * proxy, 0.0)
\end{verbatim}

When a sensitivity table is present, the price \emph{is} the difference of two measured
sensitivities. Note what disappears: the percentile rank, the $4^{-b}$ error curve
(\code{\_error\_scale}, retained only for the proxy path), and the explicit \code{num\_params}
factor --- parameter count is already inside the sum of Equation~\ref{eq:sens}. The proxy branch
survives only for modules in \code{unestimable}.

\subsection{The search}

\code{allocate\_bits(graph, scores, budget, policy, sensitivity)} runs three passes:

\begin{enumerate}
\item \textbf{Initialise at floors.} Every module starts at its role's floor width.
\item \textbf{Downgrade while over budget.} Repeatedly take the cheapest available downward move
by \code{move\_value}. If minimum-everywhere plus the hard structural floors still will not fit,
raise \code{InfeasibleTargetError} rather than silently returning something that misses the
target.
\item \textbf{Always upgrade.} Even when the floor map fits, run an upgrade pass buying the
most valuable available moves until the budget is exhausted. Discrete width steps leave slack
--- typically a few percent of the budget --- and reclaiming it is free accuracy.
\end{enumerate}

Step 3 is not an optimisation, it is a correctness fix. In the original allocator, if
$\text{budget} - \text{base} - \text{floors}$ came out negative the function returned the floor
map and the greedy loop never executed --- so at the published 3-bit setting the score had
\emph{no effect at all} and the configuration was effectively hand-written. Soft floors plus a
mandatory upgrade pass mean the score drives allocation at every target.

\subsection{Floors: soft quality, hard structure}

Two kinds of floor, and the distinction is the difference between a preference and an invariant.

\textbf{Soft quality floors} express ``this role usually wants at least $b$ bits''. They can be
breached if the budget demands it, and every breach is reported by role.

\textbf{Hard structural floors} express ``below this width the architecture stops functioning''.
No budget breaks them:

\begin{table}[H]
\centering\small
\caption{\code{STRUCTURAL\_ROLES}: hard floors, unbreakable by any budget.}
\begin{tabular}{p{0.215\textwidth}cp{0.52\textwidth}}
\toprule
Role & Floor & Why it is structural \\
\midrule
\code{MOE\_ROUTER} & 8 & $O(\text{hidden}\times\text{experts})$, often $<0.01\pct$ of the model. A
wrong argmax routes a token through the wrong 500M parameters. \\
\rowcolor{rowtint}
\code{MLA\_KV\_A}, \code{MLA\_Q\_A} & 8 & The latent bottlenecks: ${\sim}1\pct$ of the block's
parameters carrying all of its information. \\
\code{LIN\_ATTN\_A}, \code{LIN\_ATTN\_B} & 8 & ${\sim}0.03\pct$ of a gated linear-attention block
and they feed an exponential decay term. \\
\rowcolor{rowtint}
\code{SSM\_X}, \code{SSM\_DT} & 8 & Same argument: negligible size, feeds an exponential. \\
\bottomrule
\end{tabular}
\end{table}

\code{LM\_HEAD} is \emph{deliberately excluded} from that list, and the reason is worth stating
because it looks like an omission. Its 8-bit floor is a strong preference, not an invariant. On a
tied model the head \emph{is} the embedding table --- 27\pct{} of Qwen3.5-2B's parameters ---
and pinning 27\pct{} of the model at 8 bits makes every average target below ${\sim}3.9$ bits
arithmetically impossible. A hard floor there would not protect quality; it would make the tool
refuse to run.

\subsection{Fallback calibration}

Modules priced by the proxy branch and modules priced by real sensitivity live in different
units, and mixing them naively means one population is systematically cheaper and gets destroyed
first. \code{\_apply\_fallback\_scale} rescales the proxy population by a ratio of medians ---
but of \emph{rung-normalised} prices, each next-step value divided by its width span. Computing
the medians on raw step values double-counts the width-span factor and yields a scale of
$0.332$ instead of $1.000$; see defect~3 in Section~\ref{sec:defects}, where on this
architecture that error would have put 604\,M parameters at the front of the queue for the
chopping block.

\subsection{Default floors by role (excerpt)}

\begin{table}[H]
\centering\small
\caption{\code{DEFAULT\_FLOOR\_BITS}, abbreviated. \code{UNQUANTIZED\_FLOOR} keeps the tensor in
its original dtype.}
\begin{tabular}{llp{0.50\textwidth}}
\toprule
Role & Floor & Rationale \\
\midrule
\code{EMBEDDING} & 4 & Errors here hit every token. \\
\rowcolor{rowtint}
\code{LM\_HEAD} & 8 & Soft; see above. \\
\code{ATTN\_Q/K/V/O}, \code{ATTN\_QKV}, \code{ATTN\_KV} & 4 & Q/K errors compound through the
softmax. \\
\rowcolor{rowtint}
\code{ATTN\_Q\_GATE} & 4 & Fused $[Q; \text{output gate}]$; refined per row partition. \\
\code{LIN\_ATTN\_QKV/Z/OUT} & 4 & Behaves like ordinary attention. \\
\rowcolor{rowtint}
\code{LIN\_ATTN\_A/B} & 8 & Hard structural. \\
\code{LIN\_ATTN\_CONV}, \code{SSM\_CONV} & --- & Unquantized. Shape $[\text{ch},1,4]$; identical
tensors get identical answers regardless of enclosing block. \\
\rowcolor{rowtint}
\code{MLP\_GATE}, \code{MLP\_GATE\_UP} & 4 & The SwiGLU gate multiplies the up branch, so its
error is amplified rather than averaged. \\
\code{MLP\_UP}, \code{MLP\_DOWN}, \code{MLP\_FC} & 3 & The redundant bulk. \\
\rowcolor{rowtint}
\code{MOE\_ROUTER} & 8 & Hard structural. \\
\code{MOE\_EXPERT\_UP/DOWN/FC} & 2 & Each expert sees a fraction of tokens. \\
\rowcolor{rowtint}
\code{MOE\_SHARED\_*} & 3--4 & Runs for every token, so it behaves like a dense FFN. \\
\code{MLA\_Q\_A}, \code{MLA\_KV\_A} & 8 & Hard structural. \\
\rowcolor{rowtint}
\code{MLA\_Q\_B}, \code{MLA\_KV\_B} & 4 & The bulk. \\
\code{SSM\_IN}, \code{SSM\_OUT} & 4 & The bulk. \\
\rowcolor{rowtint}
\code{SSM\_X}, \code{SSM\_DT} & 8 & Hard structural. \\
\code{VISION\_PATCH\_EMBED} & 8 & A few percent of a VLM but carries all image information. \\
\rowcolor{rowtint}
\code{OTHER} & 4 & Conservative default; \code{dynquant inspect} warns and lists every
unclassified module so nothing silently drops to 2 bits. \\
\bottomrule
\end{tabular}
\end{table}

%% ============================================================ 6
\newpage
\section{The encoder}
\label{sec:encoder}

\subsection{Group-wise affine with a float offset}

Weights are grouped along the input axis in blocks of \code{group\_size} (default 128,
constrained to $\text{group\_size} \bmod 32 = 0$). Each group gets a scale and an offset, and
values reconstruct as
\begin{equation}
w \;\approx\; q \cdot \text{scale} \;+\; \text{offset},
\qquad q \in \{0, \dots, 2^b - 1\}.
\end{equation}
The offset is an \emph{unconstrained float}, not an integer zero-point. This is the single
highest-leverage decision in the format and it costs nothing at runtime: the dequantization
$q \cdot s + o$ is one FFMA either way.

The integer-zero-point convention requires the representable range to include zero, which forces
the range of any group whose values do not straddle zero to be widened all the way to it. A
worked case from the format spec: a group measured at $[0.50, 0.52]$ --- a perfectly ordinary
outcome for a post-\code{SiLU} activation path or a positively-biased weight block --- widens to
$[0.00, 0.52]$. That is a \textbf{26.0\x} range inflation, discarding \textbf{4.70 bits} of
resolution at a nominal 4-bit width. There is no reason to pay it.

Scales and offsets are rounded to the storage dtype \emph{before} the integer codes are derived,
so the codes are optimal for the parameters that will actually be read back --- not for
higher-precision parameters that then get rounded. A degenerate group (all values equal) encodes
as $\text{scale}=0$, $\text{offset}=\text{constant}$, $q=0$, which round-trips exactly and cannot
divide by zero.

\subsection{MSE-optimal clip search}

Min/max grouping is optimal for the worst case and poor for the typical case, because one outlier
stretches the grid for 127 well-behaved neighbours. The encoder therefore searches a clip ratio
per group:

\begin{verbatim}
CLIP_CANDIDATES = (1.0, 0.98, 0.96, 0.94, 0.92, 0.90, 0.85, 0.80)
\end{verbatim}

Two deliberate departures from the original formulation:

\begin{enumerate}
\item \textbf{The objective is the real codec's error.} Each candidate is pushed through
\code{QuantTensor.from\_dense} and then \code{.dequantize()}, so the winning $\alpha$ is optimal
for the bytes that reach disk --- including the dtype rounding of scale and offset, and the
degenerate-group handling. Scoring against an idealised affine map instead would pick a
different, slightly wrong winner.
\item \textbf{Selection is per group, and the ratios are returned rather than applied.} The
caller receives the chosen ratios, which makes the search reproducible and lets the sensitivity
estimator evaluate exactly the same encoding the quantizer will emit.
\end{enumerate}

\subsection{Packing}

Codes pack into \code{uint32} words with group boundaries aligned to word boundaries. At group
128 this gives $8 / 12 / 16 / 32$ words per group for $2 / 3 / 4 / 8$ bits. The alignment is what
makes an efficient kernel possible at all: a thread block loads a fixed number of words with
compile-time-constant shifts and there are no cross-lane carries. The original research code
flattened the whole tensor before 3-bit packing, so group boundaries fell mid-word and no
efficient kernel existed for that layout.

%% ============================================================ 7
\section{On-disk format}
\label{sec:format}

\subsection{Versioning}

Four numbers are versioned independently, because they change for independent reasons:

\begin{table}[H]
\centering\small
\begin{tabular}{lp{0.62\textwidth}}
\toprule
Version & Governs \\
\midrule
core version & The Python API surface. \\
\rowcolor{rowtint}
\code{CHECKPOINT\_FORMAT\_VERSION} & Layout of the weight files and manifest. A reader refuses an
unknown value rather than guessing. \\
\code{STATS\_SCHEMA\_VERSION} & The signals sidecar. v1 (research era) migrates to v2 losslessly. \\
\rowcolor{rowtint}
\code{KERNEL\_ABI\_VERSION} & The compiled extension's contract with core. Checked at import;
mismatch is a hard error with an actionable message, not a segfault later. \\
\bottomrule
\end{tabular}
\end{table}

\code{QuantTensor} stores \code{bits}, \code{group\_size}, \code{symmetric},
\code{row\_partitions}, \code{layout}, \code{shape} and \code{dtype} explicitly. None of it is
inferred. The research code stored no \code{group\_size} at all and \emph{guessed} it from
\code{scale.numel() // out\_features} with a hardcoded 128 fallback, which breaks for
\code{ndim} $\ne 2$ (Mamba \code{conv1d}, batched expert weights $[E,\text{out},\text{in}]$) and
for padded \code{in\_features}.

\subsection{Kernel chunk invariant}

The GEMV processes values in chunks sized so that a whole number of 32-bit words maps to a whole
number of values: $\text{kValues} \cdot \text{BITS} = \text{kWords} \cdot 32$. That single rule
fixes the vector load type per width:

\begin{table}[H]
\centering\small
\caption{Chunk geometry per width. The 3-bit case is the interesting one: 6 words is not a power
of two, so it loads as \code{uint2}${}\times{}$3 rather than \code{uint4}.}
\begin{tabular}{rrrl}
\toprule
Bits & Words/chunk & Values/chunk & Load type \\
\midrule
2 & 4 & 64 & \code{uint4} \\
\rowcolor{rowtint}
3 & 6 & 64 & \code{uint2} \\
4 & 4 & 32 & \code{uint4} \\
\rowcolor{rowtint}
8 & 4 & 16 & \code{uint4} \\
\bottomrule
\end{tabular}
\end{table}

Every group is word-aligned by construction (Section~\ref{sec:encoder}), so scales and offsets
are fetched once per group into shared memory and the inner loop is pure integer unpack plus
FFMA.

%% ============================================================ 8
\section{The packed runtime}
\label{sec:runtime}

\subsection{Dispatch}

\code{DynQuantLinear} holds packed \code{int32} buffers, a scale tensor and an offset tensor, and
never materialises the full-precision weight. Forward dispatch is on row count $M$:

\begin{itemize}
\item $M \le \code{gemv\_max\_rows()} = 8$: \code{dynquant::gemv\_nbit<BITS, GROUP\_SIZE>}, one
templated instantiation per width so there is no runtime branching in the inner loop.
\item $M > 8$: tile dequant into a bounded workspace, then \code{F.linear}. Correct and
throughput-competitive immediately; the fused tensor-core path is future work.
\end{itemize}

Templating on \code{BITS} rather than branching is not cosmetic. The unpack shift sequence and
the load type both depend on the width; a runtime branch inside the inner loop would cost more
than the arithmetic it guards.

\subsection{Accuracy through the kernels}

The critical property of a quantization runtime is that it does not change the answer relative to
the simulated path. Measured on the full CaseHOLD evaluation:

\begin{table}[H]
\centering\small
\caption{Simulated versus packed execution, $\approx$3.25-bit sensitivity allocation.}
\begin{tabular}{lrr}
\toprule
Path & Correct & Exact match \\
\midrule
Simulated (dequantized weights, \code{F.linear}) & 4468 / 5314 & 84.0798\pct \\
\rowcolor{rowtint}
Packed (CUDA GEMV + dequant fallback) & 4468 / 5314 & 84.0798\pct \\
\midrule
Difference & 0 & 0.0000\pct \\
\bottomrule
\end{tabular}
\end{table}

Not ``within noise'' --- identical, problem by problem. The evaluation script asserts
\code{MATCH} and fails otherwise. Worst observed relative error on any single GEMV output element
across the parity suite is $6.96\times10^{-3}$, consistent with fp16 accumulation of a
2-bit-quantized dot product.

\subsection{Memory}

\begin{table}[H]
\centering\small
\caption{VRAM, $\approx$3.25-bit allocation versus bf16. ``Manifest'' is what the allocator
predicted before anything was written.}
\begin{tabular}{lrrr}
\toprule
Measurement & Packed (GiB) & bf16 (GiB) & Reduction \\
\midrule
Manifest prediction & 0.7118 & --- & --- \\
\rowcolor{rowtint}
Walked tensors & 0.7120 & 3.5052 & 4.92\x \\
\code{torch.cuda.max\_memory\_allocated} & 0.7237 & 3.5052 & \textbf{4.84\x} \\
\rowcolor{rowtint}
\code{torch.cuda.max\_memory\_reserved} & 0.7480 & 3.5052 & 4.77\x \\
Decode peak, batch 1 & 0.882 & 3.662 & 4.15\x \\
\rowcolor{rowtint}
Decode peak, batch 32 & 5.382 & 8.164 & 1.52\x \\
\bottomrule
\end{tabular}
\end{table}

The allocator predicted 764\,317\,696 bytes; live measurement was 764\,530\,560 bytes, an
overshoot of \textbf{0.03\pct}. That is the number that makes \code{--target-size} usable: the
manifest figure is the on-disk and in-VRAM figure, with scales, offsets and packing overhead all
accounted for.

At batch 32 the reduction collapses to 1.52\x{} because activations and the KV cache dominate,
which is expected and not a weight-quantization concern.

\subsection{Speed}

\begin{table}[H]
\centering\small
\caption{Decode throughput, tokens/s. Above $M = 8$ the kernel is not used at all --- the
dequant fallback runs, and it pays for the dequant without amortising it.}
\begin{tabular}{rrrr}
\toprule
Batch & Packed & bf16 & Ratio \\
\midrule
1 & 29.9 & 33.1 & \textbf{0.90\x} \\
\rowcolor{rowtint}
4 & 116.3 & 126.2 & 0.92\x \\
8 & 235.6 & 250.9 & 0.94\x \\
\rowcolor{rowtint}
32 & 610.2 & 1012.6 & 0.60\x \\
\bottomrule
\end{tabular}
\end{table}

Decode is slower, and the reason is not the kernel. A per-step profile isolates it:

\begin{table}[H]
\centering\small
\caption{Per-step profile, batch 1 decode.}
\begin{tabular}{lrrl}
\toprule
Metric & bf16 & Packed & \\
\midrule
GPU busy time & 8.801\,ms & 7.763\,ms & $-1.038$\,ms \\
\rowcolor{rowtint}
\quad of which matmul & 3.519\,ms & 2.486\,ms & $-1.033$\,ms, \textbf{1.42\x{} faster} \\
matmul share of busy & 40.0\pct & 32.0\pct & \\
\rowcolor{rowtint}
kernel launches & 2013 & 1980 & \\
wall-clock per step & 28.4\,ms & 35.7\,ms & $+7.3$\,ms \\
\rowcolor{rowtint}
busy share of wall & 31\pct & 22\pct & \\
\bottomrule
\end{tabular}
\end{table}

The GPU does strictly less work with packed weights, and the entire $1.038$\,ms reduction in busy
time is accounted for by the $1.033$\,ms reduction in matmul time. The loss is on the host:
roughly \textbf{20--45\,\textmu s of Python and launch overhead per packed module per step},
across 187 modules. At 22\pct{} busy share the decode loop is host-bound, so the fix is CUDA
Graph capture of the decode step (planned) rather than kernel work. The 3\pct{} slowdown on full
evaluation (355.1\,s packed versus 344.6\,s bf16) is the same story diluted by
prefill-dominated batches; batch-32 prefill is 1605\,ms versus 1599\,ms, i.e.\ unchanged.

\subsection{Isolated GEMV performance}

Stripping the host overhead away and timing the kernel alone against a bf16 GEMV of the same
shape, and against a general (non-templated) quantized kernel:

\begin{table}[H]
\centering\small
\caption{Isolated GEMV, $M = 1$, A100. Bandwidth is against 1665\,GB/s measured achievable read
bandwidth (datasheet peak is 1935\,GB/s). ``Gain'' is speedup over the general kernel.}
\begin{tabular}{llrrrr}
\toprule
Module & Shape & Bits & Speedup vs bf16 & \pct{} of 1665\,GB/s & Gain \\
\midrule
\code{out\_proj}    & $2048\times2048$ & 2 & 1.23\x & 10\pct & 1.52\x \\
                    &                  & 3 & 1.23\x & 15\pct & 1.55\x \\
                    &                  & 4 & 1.18\x & 18\pct & 1.78\x \\
                    &                  & 8 & 1.09\x & 33\pct & 1.94\x \\
\rowcolor{rowtint}
\code{in\_proj\_qkv}, \code{gate\_proj} & $6144\times2048$ & 2 & 1.26\x & 20\pct & 1.32\x \\
\rowcolor{rowtint}
                    &                  & 3 & 1.24\x & 29\pct & 1.53\x \\
\rowcolor{rowtint}
                    &                  & 4 & 1.21\x & 37\pct & 1.47\x \\
\rowcolor{rowtint}
                    &                  & 8 & 1.16\x & 68\pct & 1.63\x \\
\code{down\_proj}   & $2048\times6144$ & 2 & 1.51\x & 14\pct & 1.81\x \\
                    &                  & 3 & 1.50\x & 20\pct & 1.85\x \\
                    &                  & 4 & 1.40\x & 25\pct & 2.09\x \\
                    &                  & 8 & 1.29\x & 44\pct & 2.36\x \\
\rowcolor{rowtint}
\code{embed}/\code{lm\_head} & $248320\times2048$ & 2 & \textbf{2.56\x} & 36\pct & 1.36\x \\
\rowcolor{rowtint}
                    &                  & 3 & \textbf{2.53\x} & 52\pct & 1.96\x \\
\rowcolor{rowtint}
                    &                  & 4 & \textbf{2.51\x} & 67\pct & 1.61\x \\
\rowcolor{rowtint}
                    &                  & 8 & 1.89\x & 98\pct & 1.44\x \\
\bottomrule
\end{tabular}
\end{table}

The kernel is faster than bf16 at every width and shape, by 1.09\x{} to 2.56\x. It also beats a
general quantized kernel by 1.32\x{} to 2.36\x, which is what the \code{<BITS, GROUP\_SIZE>}
templating buys.

But the bandwidth gate --- $\ge$70\pct{} of achievable --- is met only at 8 bits on the large
tensor (98\pct). At 4, 3 and 2 bits the best figures are 67\pct, 52\pct{} and 36\pct.
Section~\ref{sec:notestablished} explains why this is arithmetic, not a bug.

\subsection{Verification}

\begin{table}[H]
\centering\small
\caption{Verification of the shipping binary.}
\begin{tabular}{lr}
\toprule
Check & Result \\
\midrule
Kernel parity vs torch oracle, all (bits, group, shape, $M$) & 417 passed \\
\rowcolor{rowtint}
Full test suite & 1169 passed, 34 skipped, 0 failed \\
\code{compute-sanitizer --tool memcheck} & 417 passed, 0 errors \\
\rowcolor{rowtint}
\code{compute-sanitizer --tool initcheck} & 417 passed, 0 errors \\
\code{compute-sanitizer --tool racecheck} & 417 passed, \textbf{0 hazards} (0 errors, 0 warnings) \\
\rowcolor{rowtint}
\code{compute-sanitizer --tool synccheck} & 417 passed, 0 errors \\
Register usage & 64 (72 at 3 bits), \code{STACK:0} --- no spills \\
\rowcolor{rowtint}
Determinism & Same input $\to$ bit-identical output across runs \\
\bottomrule
\end{tabular}
\end{table}

These ran against the exact binary that produced the performance numbers. The build was
re-verified from a clean rebuild of the shipping source after late edits to \code{gemv.cu}: the
rebuild is codegen-identical --- 1384 SASS instructions in the 2-bit $M=1$ kernel with the same
\code{FFMA} 264 / \code{I2F.U16} 256 / \code{LOP3} 244 / \code{SHF} 240 histogram, the same 64
registers and \code{STACK:0} --- and re-measures the \code{embed}/\code{lm\_head} row at
2.55/2.54/2.51/1.90\x{} and 36/52/67/99\pct, the same numbers inside run-to-run noise. The 417
parity tests are what actually validates the host-side grid-expression change, since an undersized
grid leaves output rows untouched rather than faulting.

%% ============================================================ 9
\newpage
\section{Experimental setup}

\begin{table}[H]
\centering\small
\caption{Experimental configuration.}
\begin{tabular}{ll}
\toprule
Model & Qwen3.5-2B-Base, \code{Qwen3\_5ForCausalLM} \\
\rowcolor{rowtint}
Architecture & 28 layers, hybrid gated-linear / full attention, tied embeddings, \\
\rowcolor{rowtint}
 & fused \code{q\_proj} $=[Q;\ \text{output gate}]$ (\code{attn\_output\_gate: true}) \\
Parameters & 1.8821\,B total; 1.8817\,B across 187 quantizable modules \\
\rowcolor{rowtint}
Hardware & NVIDIA A100 80\,GB PCIe, 108 SMs, 40\,MB L2 \\
Software & CUDA 13.0.88, PyTorch 2.13.0+cu130, Transformers 5.14.1 \\
\rowcolor{rowtint}
Task & CaseHOLD --- five-way legal holding selection, exact match, $n = 5314$ \\
Calibration & CaseHOLD \code{validation}, 512 rows, 204\,693 tokens \\
\rowcolor{rowtint}
Group size & 128 throughout \\
Widths & $\{2, 3, 4, 8\}$ \\
\bottomrule
\end{tabular}
\end{table}

\subsection{Choosing the task: the headroom screen}

Before spending GPU-hours, four candidate datasets were screened for whether fine-tuning could
move the base model at all. The rule: a dataset is only useful if the base model is well above
chance (so the task is learnable) \emph{and} well below the supervised reference (so there is
headroom for quantization to then lose).

\begin{table}[H]
\centering\small
\caption{Headroom screen. Base accuracy measured; supervised reference from the literature.}
\begin{tabular}{lrrrl}
\toprule
Dataset & Base & Chance & Supervised ref. & Verdict \\
\midrule
\rowcolor{rowtint}
CaseHOLD & 34.3\pct & 20\pct & ${\sim}69\pct$ & \textbf{Pass} --- used \\
sql-create-context & 54.0\pct & 0\pct & ${\sim}80\pct$ & Pass \\
MedMCQA & 47.7\pct & 25\pct & ${\sim}45$--$50\pct$ & Reject --- base already at reference \\
GSM8K & 66.1\pct & 0\pct & 65.1\pct measured & Reject --- base \emph{above} reference \\
\bottomrule
\end{tabular}
\end{table}

GSM8K is the instructive failure. The base model already scores 66.11\pct; fine-tuning moves it
to 65.13\pct, a change of $-0.99$\pt{} at $p = 0.48$. There is no gain, so there is nothing for
quantization to retain, and every quantization arm on that task measures only degradation from a
ceiling. The dataset was run anyway (Section~\ref{sec:gsm8k}) and its results are reported,
because a flat arm is itself a finding --- but the design lesson is to measure base accuracy
\emph{before} committing to a fine-tune.

%% ============================================================ 10
\section{Results}

\subsection{Accuracy}

\begin{table}[H]
\centering\small
\caption{CaseHOLD exact match, $n = 5314$. The base model produced 7 unparseable outputs, counted
as wrong.}
\begin{tabular}{lrrr}
\toprule
Arm & Correct & Exact match & Stored bits \\
\midrule
Base, fp16 & 1867 & 35.13\pct & 16.000 \\
\rowcolor{rowtint}
Fine-tuned, fp16 & 4702 & \textbf{88.48\pct} & 16.000 \\
DynQuant $\approx$4.25\,b & 4596 & 86.49\pct & 4.249 \\
Uniform 4\,b & 4603 & 86.62\pct & 4.255 \\
\rowcolor{rowtint}
\textbf{DynQuant $\approx$3.25\,b, sensitivity} & \textbf{4468} & \textbf{84.08\pct} & 3.249 \\
Uniform 3\,b & 3921 & 73.79\pct & 3.256 \\
DynQuant $\approx$3.25\,b, rank product & 3813 & 71.75\pct & 3.249 \\
\midrule
Guessing & --- & 20.00\pct & --- \\
\bottomrule
\end{tabular}
\end{table}

\subsection{Paired significance}

All arms were evaluated on the same 5314 problems and per-problem correctness was recorded, so
every comparison below is an \emph{exact} paired McNemar test, not a two-proportion
approximation. Confidence intervals are 95\pct{} on the paired difference.

\begin{table}[H]
\centering\small
\caption{Exact paired McNemar, same 5314 problems throughout.}
\begin{tabular}{lrrr}
\toprule
Comparison & $\Delta$ (pts) & 95\pct{} CI & $p$ \\
\midrule
\rowcolor{rowtint}
Fine-tune vs base & $+53.35$ & $[51.87,\ 54.83]$ & $<10^{-300}$ \\
\midrule
Cost of 4-bit quantization & $-1.99$ & $[-2.68,\ -1.31]$ & $1.3\times10^{-8}$ \\
\rowcolor{rowtint}
Allocation vs uniform @ 4\,b & $-0.13$ & $[-0.51,\ +0.25]$ & $0.56$ \emph{(not separated)} \\
\midrule
Cost of 3-bit quantization (sensitivity) & $-4.40$ & $[-5.23,\ -3.57]$ & $4.8\times10^{-26}$ \\
\rowcolor{rowtint}
Cost of 3-bit quantization (rank product) & $-16.73$ & --- & $1.1\times10^{-155}$ \\
\midrule
\rowcolor{rowtint}
\textbf{Sensitivity vs uniform @ 3\,b} & $\mathbf{+10.29}$ & $[9.20,\ 11.39]$ & $1.3\times10^{-81}$ \\
Rank product vs uniform @ 3\,b & $\mathbf{-2.03}$ & $[-3.01,\ -1.05]$ & $5.8\times10^{-5}$ \\
\rowcolor{rowtint}
\textbf{Sensitivity vs rank product @ 3\,b} & $\mathbf{+12.33}$ & $[11.13,\ 13.53]$ & $1.7\times10^{-96}$ \\
\bottomrule
\end{tabular}
\end{table}

Reading the table in order:

\begin{itemize}
\item Fine-tuning works, enormously, and that gain is the thing being protected.
\item At 4.25 bits, quantization costs $2$\pt{} and \emph{allocation does not matter} --- the
$-0.13$\pt{} difference from uniform is not distinguishable from zero. When the budget is loose
enough that nothing has to be hurt, there is nothing to allocate.
\item At 3.25 bits, allocation is the whole game. Sensitivity-driven allocation costs $4.40$\pt{}
where uniform costs $14.69$ and the published rank product costs $16.73$.
\item The rank product is \emph{significantly worse than uniform} at this budget. That is a
stronger statement than ``it does not help''. A score that loses to no score is a score pointing
the wrong way.
\end{itemize}

\subsection{Retention of the fine-tuning gain}

\begin{table}[H]
\centering\small
\caption{Share of the $+53.35$\pt{} fine-tuning gain surviving each arm.}
\begin{tabular}{lr}
\toprule
Arm & Retention \\
\midrule
fp16 & 100\pct \\
\rowcolor{rowtint}
$\approx$4.25\,b & 96\pct \\
\textbf{$\approx$3.25\,b, sensitivity} & \textbf{92\pct} \\
$\approx$3.25\,b, rank product & 69\pct \\
\bottomrule
\end{tabular}
\end{table}

This is the framing that matters for a practitioner. You spent a fine-tune to gain 53 points. At
a 4.9\x{} compression ratio you keep 92\pct{} of it with the shipped estimator, or 69\pct{} with
the published one --- and the difference is entirely in which tensors got which widths, at
identical total size.

\subsection{The answer key, and what correlates with it}
\label{sec:answerkey}

To evaluate a sensitivity estimator you need ground truth: the actual loss damage attributable to
each module. Two such keys were built, and they answer subtly different questions. Every
correlation in the literature and in this paper should be labelled with which one it used,
because the numbers differ.

\begin{table}[H]
\centering\small
\caption{The two ground-truth constructions. Neither is wrong; they measure different things.}
\begin{tabular}{p{0.14\textwidth}p{0.38\textwidth}p{0.38\textwidth}}
\toprule
& \textbf{Key A} --- \code{marginal\_3to4.json} & \textbf{Key B} --- \code{sensitivity\_3bit.json} \\
\midrule
Construction & Whole model at 3 bits; each module \emph{restored} to 4 bits, one at a time &
Whole model at fp16; each module quantized \emph{alone} to 3 bits \\
\rowcolor{rowtint}
Target & Signed \code{gain} in loss & $|\Delta\text{loss}|$, disturbance magnitude \\
Asks & ``what does one more bit here buy, in context?'' & ``how much does this module dislike
3 bits, in isolation?'' \\
\bottomrule
\end{tabular}
\end{table}

Key A is the one that matches what the allocator actually does, so it is the primary key.

\begin{table}[H]
\centering\small
\caption{Spearman $\rho$ against \textbf{Key A} (signed 3$\to$4 gain). ``Within-role'' is the mean
over 12 roles, with the count of roles in which the sign is positive.}
\begin{tabular}{lrr}
\toprule
Candidate signal & Global $\rho$ & Within-role mean ($n$ positive / 12) \\
\midrule
Plasticity (gradient variance) & $+0.3131$ & $+0.266$ \ (10/12) \\
\rowcolor{rowtint}
\textbf{Fisher / Gauss--Newton (Eq.~\ref{eq:sens})} & $\mathbf{+0.2951}$ & $\mathbf{+0.281}$ \ (11/12) \\
\code{num\_params} & $+0.1492$ & --- \\
\rowcolor{rowtint}
Weight reconstruction error only & $+0.0645$ & --- \\
Gram (activation second moment only) & $+0.0099$ & $-0.151$ \ (4/12) \\
\rowcolor{rowtint}
\textbf{Shipped rank product (Eq.~\ref{eq:oldscore})} & $\mathbf{+0.0040}$ & $-0.001$ \ (5/12) \\
Saliency ($\operatorname{EMA}\lVert X\rVert$) & $\mathbf{-0.2059}$ & $-0.214$ \ (2/12) \\
\bottomrule
\end{tabular}
\end{table}

\begin{table}[H]
\centering\small
\caption{Within-role Spearman $\rho$ against \textbf{Key B} ($|\Delta\text{loss}|$ from isolated
3-bit quantization), mean over 12 roles.}
\begin{tabular}{lrr}
\toprule
Candidate signal & Within-role mean & Positive roles \\
\midrule
\rowcolor{rowtint}
\textbf{Gauss--Newton (Eq.~\ref{eq:sens})} & $\mathbf{+0.521}$ & 11/12 \\
Plasticity & $+0.491$ & 11/12 \\
\rowcolor{rowtint}
Shipped rank product & $+0.231$ & 9/12 \\
Weight-only & $-0.227$ & 4/12 \\
\rowcolor{rowtint}
Saliency & $-0.301$ & 2/12 \\
Eq.~\ref{eq:sens} with $\mathbb{E}[\delta^2]$ removed & $\mathbf{-0.338}$ & 1/12 \\
\bottomrule
\end{tabular}
\end{table}

The two keys agree on everything that matters, which is why both are reported:

\begin{itemize}
\item \textbf{Saliency anti-correlates}, on both keys, in 10 of 12 roles. Busy tensors are robust
tensors on this architecture. The published formula multiplies by this.
\item \textbf{The gradient factor is the load-bearing one.} Removing $\mathbb{E}[\delta^2]$ from
Equation~\ref{eq:sens} does not merely weaken it --- it flips the sign, from $+0.521$ to
$-0.338$, positive in 1 role out of 12. Activation statistics alone are worse than nothing.
\item \textbf{The rank product is approximately noise} on the key that matches the allocator's
decision ($+0.0040$ globally, $-0.001$ within role, positive in 5 of 12 roles --- a coin flip).
\item \textbf{Plasticity and Gauss--Newton are the two real signals}, and they are close.
Gauss--Newton wins on within-role consistency (11/12 on both keys) and, decisively, in the
end-to-end shootout below.
\end{itemize}

For scale: fp16 loss on the calibration set is $0.1049$; uniform 3-bit loss is $0.1514$. So
$0.0466$ of loss is the entire recoverable budget that any allocation policy is fighting over.

\subsection{Policy shootout}
\label{sec:shootout}

Each candidate policy was run end to end --- allocate, quantize, evaluate loss --- at two tight
budgets. The metric is the share of uniform-3-bit damage recovered, where 100\pct{} means
matching an oracle that knows the true per-module damage. Values above 100\pct{} are possible
because the oracle is itself greedy on a first-order key.

\begin{table}[H]
\centering\small
\caption{Share of uniform-3-bit damage recovered. Higher is better; negative means worse than
uniform.}
\begin{tabular}{lrr}
\toprule
Policy & @ 3.125 bits & @ 3.250 bits \\
\midrule
Oracle (knows the true key) & 84.87\pct & 104.15\pct \\
\midrule
\code{fisher\_move} & 77.21\pct & \textbf{105.67\pct} \\
\rowcolor{rowtint}
\textbf{\code{fisher\_diff} --- ships} & \textbf{85.54\pct} & 95.76\pct \\
Plasticity & 86.81\pct & 81.58\pct \\
\midrule
Shipped rank product (Eq.~\ref{eq:oldscore}) & 28.49\pct & 85.17\pct \\
\rowcolor{rowtint}
Weight reconstruction error only & $\mathbf{-29.46\pct}$ & 48.51\pct \\
\code{num\_params} & 52.01\pct & 26.40\pct \\
\midrule
\rowcolor{rowtint}
Role-aware floors, no signal at all & 65.58\pct & 95.08\pct \\
Random, 5 seeds & $-6.45$ to $45.39\pct$ & $-5.86$ to $59.15\pct$ \\
\bottomrule
\end{tabular}
\end{table}

The shipping policy \code{fisher\_diff} --- Equation~\ref{eq:sens} with difference pricing --- is
the only candidate above 85\pct{} at \emph{both} budgets. Three other readings deserve
attention:

\begin{itemize}
\item \textbf{The tighter budget is the discriminating one.} At 3.250 bits, five policies land in
the 85--106\pct{} band and even ``role-aware floors with no signal'' reaches 95.08\pct. At 3.125
bits the field spreads from $-29$\pct{} to $+87$\pct. Evaluating an allocator only at a loose
budget will fail to distinguish it from a lookup table.
\item \textbf{Calibration-free weight error is actively harmful} at the tight budget
($-29.46$\pct, worse than every random seed). Reconstruction error without the two moment vectors
is not a usable importance proxy.
\item \textbf{Random allocation can look respectable.} One seed reached $+59.15$\pct{} at
3.250 bits. Any published allocator result at a loose budget without a random baseline should be
treated as unvalidated.
\end{itemize}

\subsection{What the allocation actually does differently}

At an identical 3.249 stored bits, here is where the sensitivity estimator puts the budget
relative to the rank product:

\begin{table}[H]
\centering\small
\caption{Average bits per role at $\approx$3.25 bits. Both columns are the same total size.}
\begin{tabular}{lrrr}
\toprule
Role & Sensitivity & Rank product & $\Delta$ \\
\midrule
\rowcolor{rowtint}
\code{k\_proj} & \textbf{4.00} & 3.17 & $+0.83$ \\
\rowcolor{rowtint}
\code{v\_proj} & \textbf{4.00} & 3.50 & $+0.50$ \\
\code{q\_proj} (fused $Q$ + gate) & 3.33 & 3.00 & $+0.33$ \\
\rowcolor{rowtint}
\code{lin\_attn.out\_proj} & 3.33 & 3.11 & $+0.22$ \\
\code{in\_proj\_qkv} & 3.33 & 3.17 & $+0.16$ \\
\rowcolor{rowtint}
\code{o\_proj} & 3.33 & 3.17 & $+0.16$ \\
\code{in\_proj\_z} & 3.17 & 3.11 & $+0.06$ \\
\rowcolor{rowtint}
\code{embed\_tokens} & 3.00 & 3.00 & $0$ \\
\code{in\_proj\_a}, \code{in\_proj\_b} & 8.00 & 8.00 & $0$ (hard floor) \\
\rowcolor{rowtint}
\code{gate\_proj} & 3.12 & 3.12 & $0$ \\
\code{down\_proj} & \textbf{2.67} & 2.79 & $-0.12$ \\
\rowcolor{rowtint}
\code{up\_proj} & \textbf{2.67} & 2.83 & $-0.16$ \\
\bottomrule
\end{tabular}
\end{table}

The trade is legible: \textbf{take bits out of the FFN bulk and put them into K and V}. K and V
projections go all the way to 4 bits --- their errors compound through the softmax and then
through every subsequent token's attention --- paid for by pushing \code{up\_proj} and
\code{down\_proj} below 2.7 bits average, where the FFN's redundancy absorbs it.

The width histogram is the surprise:

\begin{table}[H]
\centering\small
\caption{Parameters at each extreme width. The \emph{better} map is the \emph{more} aggressive
one.}
\begin{tabular}{lrrrr}
\toprule
Width & Sensitivity & \pct{} of model & Rank product & \pct{} of model \\
\midrule
2-bit & 440.4\,M & 23.4\pct & 220.2\,M & 11.7\pct \\
\rowcolor{rowtint}
4-bit & 427.8\,M & 22.7\pct & 207.6\,M & 11.0\pct \\
\bottomrule
\end{tabular}
\end{table}

The winning allocation puts \emph{twice} as many parameters at 2 bits and twice as many at 4
bits: it is more polarised in both directions. This refutes a reading that came out of earlier
experiments --- that 2-bit round-to-nearest is simply unusable and should be avoided. It is
unusable in the wrong places. Put 23\pct{} of the model at 2 bits and spend the savings on K/V,
and you gain 12.33 points over the map that hedged.

\subsection{The GSM8K arm}
\label{sec:gsm8k}

Reported for completeness, and because it is a clean negative result.

\begin{table}[H]
\centering\small
\caption{GSM8K, $n = 1319$, exact match.}
\begin{tabular}{lrl}
\toprule
Arm & Exact match & \\
\midrule
Base, fp16 & 66.11\pct & \\
\rowcolor{rowtint}
Fine-tuned, fp16 & 65.13\pct & $-0.99$\pt, $p = 0.48$ --- no gain \\
$\approx$4.25\,b DynQuant & 58.15\pct & (uniform 4\,b: 56.79\pct) \\
\rowcolor{rowtint}
$\approx$3.25\,b DynQuant & 13.57\pct & (uniform 3\,b: 21.46\pct); allocation $-7.88$\pt, $p = 1.1\times10^{-10}$ \\
\bottomrule
\end{tabular}
\end{table}

Two things to take from this. First, the fine-tune produced no gain --- the base model was
already at the supervised reference --- so the moments were harvested from a model with nothing
to say about this task, and the allocation they produced is worse than uniform. Signals from a
converged, gainless fine-tune are not informative signals. Second, chain-of-thought arithmetic
degrades far more steeply under 3-bit quantization than five-way classification does: 65\pct{}
$\to$ 14\pct{} versus 88\pct{} $\to$ 84\pct. Multi-step generative reasoning has no error
tolerance, because a single wrong token invalidates the rest of the chain.

The 3.25-bit \emph{sensitivity} arm was not run on GSM8K, so the fair comparison --- does the new
estimator also fail on a gainless fine-tune, or does it recover? --- remains open.

%% ============================================================ 11
\newpage
\section{What is not established}
\label{sec:notestablished}

This section exists because everything above is a measurement and measurements have scope.

\subsection{Decode is slower, and the bandwidth gate is missed}

At batch 1, packed decode runs at \textbf{0.90\x} bf16 throughput. The kernel itself is faster
(Section~\ref{sec:runtime}); the loss is 20--45\,\textmu s of host overhead per module per step
over 187 modules, with the decode loop only 22\pct{} GPU-busy. CUDA Graph capture is the fix and
is not yet implemented.

The bandwidth gate of $\ge$70\pct{} achievable is met only at 8 bits on the largest tensor. The
reason is an issue-rate floor, not a memory-system problem. Analysis on sm\_80, where
FFMA/FADD/integer-add/\code{LOP3}/\code{SHF} all issue at 64 per SM per clock and conversions
from 8- and 16-bit integers at 64 while all other conversions issue at 16:

\begin{itemize}
\item The instruction mix of the 2-bit kernel gives an issue-rate floor of ${\sim}141$\,\textmu s
against ${\sim}237$\,\textmu s measured --- roughly 60\pct{} issue efficiency. The kernel is
\emph{instruction}-bound, and at low bit widths there are simply more unpack instructions per
byte of memory traffic.
\item At $K = 2048$ with 64 values per chunk, each warp executes exactly \emph{one} main-loop
iteration. There is no second iteration to overlap with, so no intra-warp memory-level
parallelism is available to hide the load latency. A rows-per-warp sweep confirms the ceiling is
structural: 2 rows gives 34/50/65/98\pct, 4 rows gives 36/52/67/98--99\pct{} (shipped), 8 rows
gives 36/49/64/97\pct{} and spills at $M \ge 4$.
\item The fp32 magic-number dequantization trick used by Marlin and AWQ, which would eliminate
the integer$\to$float conversions entirely, is \emph{not available} in this format. It requires
$\text{offset}' = -(2^{23} + z)\cdot\text{scale}$, which at 2 bits produces ${\sim}17\pct$ error.
The escape hatch those kernels use exists only in fp16, and this runtime accumulates in fp32.
\end{itemize}

So the honest statement is: the memory saving is real and verified (4.84\x); the speed claim is
not. The kernel wins in isolation and loses in the loop.

\subsection{Scope of the accuracy result}

One model (Qwen3.5-2B-Base), one task (CaseHOLD), two budgets (3.25 and 4.25 bits), one
group size (128). The $+10.29$\pt{} result is a large, robustly significant effect on that
configuration and nothing has yet shown it generalises. Specifically not established:

\begin{itemize}
\item \textbf{Larger models.} 2B is small. The redundancy structure that makes 23\pct{} of the
model tolerable at 2 bits may be quite different at 14B or 70B.
\item \textbf{MoE, MLA and SSM architectures.} The roles, floors and structural invariants are
implemented and unit-tested against tiny synthetic configs, but no MoE or MLA model has been
quantized and evaluated end to end.
\item \textbf{Generative reasoning tasks.} The GSM8K arm suggests 3-bit is simply too aggressive
for chain-of-thought, independent of allocation.
\item \textbf{Task specificity of the moments.} The central premise --- that signals harvested
during \emph{this} fine-tune produce an allocation specialised to \emph{this} task --- is
untested. At 4.25 bits the two tasks produced 177 of 187 identical widths, which is what a
task-\emph{in}dependent allocation looks like. The isolating experiment (CaseHOLD moments versus
GSM8K moments on one fixed model at one fixed budget) has not been run.
\end{itemize}

\subsection{Limits of the estimator itself}

\begin{itemize}
\item \textbf{First-order in the module.} Equation~\ref{eq:sens} prices each module as if it were
the only one quantized. Summed over an allocation it overshoots measured loss by ${\sim}2\times$.
This is why \code{absolute} is \code{False} and why the table is used as an ordering.
\item \textbf{Diagonal Gauss--Newton.} Off-diagonal curvature within a module is discarded, as is
all cross-module curvature.
\item \textbf{Post-hoc, not streamed.} Moments are written at the end of fine-tuning and read
back. There is no online allocation that adapts during training.
\item \textbf{Sampled every 16 steps.} Chosen for cost, not accuracy. The sensitivity of the
result to that interval has not been swept.
\end{itemize}

\subsection{Cost}

Packing 187 modules runs on CPU in 447\,s. Full CaseHOLD evaluation takes 355.1\,s packed versus
344.6\,s bf16.

%% ============================================================ 12
\section{Defects found and fixed}
\label{sec:defects}

Four defects are documented here because each one was silent --- it produced plausible numbers
rather than an error --- and because each says something about where this class of system fails.

\subsection{Every role's least important module was free to destroy}

\code{percentile\_ranks} mapped rank to $(\text{rank} - 1)/(n - 1)$. The minimum-ranked module in
every role therefore scored \emph{exactly} 0. Since the price of a downward move is proportional
to the score, a zero score means zero damage, so the allocator would take that module to the
minimum width first, every time, at no cost --- regardless of what the module was.

The fix is the Hazen plotting position, $(\text{rank} - 0.5)/n$, which never reaches 0 or 1.

The instructive part is how it survived: \textbf{the scorer had no test file at all}. A
score-to-rank transformation is exactly the kind of pure function that is trivially testable ---
ties, $n = 1$, all-equal inputs, NaN --- and the absence of the test was the defect that let the
other defect live.

\subsection{The paired test was unrecoverable after the GPU-hours were spent}

The evaluation harness stored summary counts per arm: correct, total, accuracy. That is enough
for a two-proportion test and \emph{not} enough for a paired one, and the arms differ by a couple
of points on 5314 problems, where pairing is worth roughly an order of magnitude in statistical
power.

By the time this was noticed, the GPU-hours were spent and the per-problem outcomes were gone.
Every arm had to be re-run.

The fix is structural: the results schema now \emph{requires} an unsampled per-problem \code{hits}
boolean array, and \code{eval/compare.py} computes exact McNemar from it. Every $p$-value in
Section~\ref{sec:answerkey} exists because of that change.

\subsection{Fallback calibration double-counted the width span}

\code{\_apply\_fallback\_scale} equalises the proxy-priced and sensitivity-priced populations by a
ratio of medians. The medians were taken over \emph{raw} next-step values --- but a next-step value
already contains the width span of that step. Dividing by the span (rung-normalising) before
taking the median is required; without it the computed scale came out $0.332$ instead of
$1.000$, a factor-of-three systematic discount on one population.

The consequence is not a small accuracy wobble. A three-fold underpricing means whichever
population sat at the lower floor becomes the cheapest thing in the queue and gets cut first ---
on this architecture, 604\,M parameters.

\subsection{A tied embedding table arrived on the GPU twice}

\code{nn.Module.\_apply} invokes its conversion function once per registered buffer and memoizes
nothing. A tied vocabulary table registered as a buffer on \emph{both} \code{embed\_tokens} and
\code{lm\_head} is therefore converted twice and \emph{two independent copies} land on the
device. Measured directly: 8192 bytes moved for one 4096-byte tensor.

On Qwen3.5-2B the vocabulary table is 27\pct{} of the model, so the headline memory number would
have been \textbf{0.9161 GiB instead of 0.7118 GiB} --- a 27\pct{} inflation of the central
claim of the runtime work, reported in perfect good faith, with the manifest and the live
measurement agreeing with each other because both were wrong the same way.

The fix is structural rather than a deduplication pass: the tied module registers no buffer at
all and reaches the owner's storage through a \code{holder} property. There is one buffer, so
there is nothing to deduplicate.

\subsection{The pattern}

All four are silent failures that yield plausible output: a zero that means ``free'' instead of
``unranked''; a summary statistic that discards the pairing; a calibration constant off by 3\x;
a memory measurement inflated by 27\pct{} with two independent sources agreeing. None would have
been caught by inspection of the results. Each was caught by asking a specific quantitative
question --- what \emph{exactly} does the minimum-ranked module score, what \emph{exactly} does
the allocator predict versus what is live --- and following the answer.

%% ============================================================ 13
\section{Using it}

\subsection{Collecting signals}

\begin{verbatim}
from dynquant import DynQuantCallback

trainer = Trainer(
    model=model,
    ...,
    callbacks=[DynQuantCallback(output_dir="stats/")],
)
trainer.train()
\end{verbatim}

\code{SFTTrainer} from TRL takes the same callback. For a hand-written training loop:

\begin{verbatim}
from dynquant import track_signals

with track_signals(model, out="stats/"):
    for batch in loader:
        loss = model(**batch).loss
        loss.backward()
        optimizer.step(); optimizer.zero_grad()
\end{verbatim}

Under DDP / FSDP / DeepSpeed the accumulators are reduced across ranks before the sidecar is
written.

\subsection{Inspecting, allocating, quantizing}

The sidecar written above drives everything downstream. \code{--moments} is what selects the
sensitivity estimator of Section~\ref{sec:sens}; without it the allocator falls back to
plasticity ranks, which is the right behaviour for a run whose moments were never collected and
the wrong thing to do by accident.

\begin{verbatim}
dynquant inspect  out/merged --stats stats/ --moments stats/ \
                  --target 3.25 --uniform 3 4 --save-map maps/
dynquant quantize out/merged --map maps/ -o out/q3
dynquant eval     out/merged --task casehold --out runs/bf16.json
dynquant eval     out/q3     --task casehold --out runs/q3.json \
                  --compare runs/bf16.json
dynquant bench    --model out/merged
\end{verbatim}

Budgets may be given as \code{--target} (average bits), \code{--target-size 6.5GiB}, or
\code{--target-ratio 0.2}. All three are accounted honestly --- scales, offsets and packing
overhead included --- so the manifest number is the on-disk number, verified here to 0.03\pct.

Three properties of that sequence are deliberate. Passing \code{--save-map} to \code{inspect} and
\code{--map} to \code{quantize} means the map that was reviewed is the map that gets applied,
with no second allocation in between that could differ from the one that was looked at.
\code{eval --compare} takes the \emph{record} written by an earlier \code{--out}, not another
model, because the test is paired on per-problem outcomes (Section~\ref{sec:answerkey}) and it
refuses to run across a differing task, split, shot count, seed or limit. And \code{bench}
measures rectangles rather than a model end to end: decode is memory-bound, so the figure that
means something is the fraction of \emph{this card's measured} achievable bandwidth that the
packed GEMV reaches, and a datasheet peak is the wrong denominator.

\code{dynquant inspect} prints the role / parameters / sensitivity / assigned-bits table and names
every module that fell through to \code{OTHER}, so an unrecognised architecture is loud rather
than quietly conservative. It also reports within-role concordance of assigned width with score,
which is the diagnostic that caught the legacy allocator ignoring its scores outright: inverting
every one of them changed zero of Qwen3-14B's 282 assigned widths, and nothing the allocator
printed said so.

\paragraph{What \code{quantize} writes today.} The packed-checkpoint writer is not yet
implemented, so \code{-o} produces an ordinary \code{transformers} directory in which every value
has been encoded to its allocated width and decoded back. Those are exactly the values a packed
checkpoint would hold, so accuracy measured on that directory is the quantized model's accuracy;
its size is the compute dtype's size, and the command prints the packed size alongside it on
every run. The configuration in which the weights really are packed in VRAM --- and therefore the
only one from which a memory or throughput figure is meaningful --- is \code{dynquant eval
out/merged --map maps/}, which swaps the modules onto the packed runtime in place. Every memory
number in this report was measured that way.

\subsection{Loading}

Today a quantized model is reconstituted in one of two ways, and which one is right depends on
what is being measured. For accuracy, the directory written by \code{quantize} loads with stock
\code{transformers} and no DynQuant installed at all. For memory or throughput, the packed
runtime is installed over a loaded model in process:

\begin{verbatim}
from transformers import AutoModelForCausalLM
from dynquant.runtime.linear import pack_model

model = AutoModelForCausalLM.from_pretrained("out/merged", ...)
report = pack_model(model, bits, group_size=128)   # weights now packed in VRAM
\end{verbatim}

\code{DYNQUANT\_BACKEND=cuda|triton|torch} forces a backend; all three agree numerically, and the
torch path is the correctness oracle that the compiled kernels are tested against.

\paragraph{Not yet implemented.} The intended end state is a checkpoint that is packed on disk
and loaded through \code{transformers} directly:

\begin{verbatim}
model = AutoModelForCausalLM.from_pretrained("my-org/qwen35-2b-dynquant-3bit")
\end{verbatim}

with \code{DynQuantHfQuantizer} registered into \code{transformers.quantizers.auto} on import and
swapping Linears \emph{during} load, so peak host memory never reaches the full-precision size.
That requires the packed-checkpoint writer and the streaming load path, neither of which is
written. It is stated here as design, not as behaviour: the format
(Section~\ref{sec:format}) was specified to admit it, and the runtime modules it would install
are the ones measured in Section~\ref{sec:runtime}.

\subsection{Reproducing the published numbers}

The superseded scorer is preserved rather than deleted, because a claim that the replacement is
better is only checkable if the thing it replaced still runs. It is reachable from the Python
API:

\begin{verbatim}
from dynquant.score.importance import ScoreConfig, score_modules

report = score_modules(graph, stats, ScoreConfig(combine="rank_product"))
\end{verbatim}

\code{combine = "rank\_product"} is Equation~\ref{eq:oldscore}, ordinal ranks and all; it is the
ordering that scores $+0.231$ in Section~\ref{sec:answerkey} against the replacement's $+0.521$.

\paragraph{Not yet implemented.} A single \code{--preset paper-3.15} flag bundling all three
compatibility switches --- \code{combine = "rank\_product"}, \code{soft\_floors = False} (an
over-budget floor map returned unchanged, so the greedy loop never runs, which is the defect of
Section~\ref{sec:alloc}) and \code{use\_coherence = True} (the EMA term present in the research
code though absent from the published formula) --- is specified but not written, and the golden
test that would pin it does not exist. The correct implementation is the default; compatibility
is opt-in and, for the moment, opt-in one flag at a time.

%% ============================================================ 14
\section{Conclusions}

\textbf{On the method.} Replacing an ordinal rank product with a cardinal, loss-unit
Gauss--Newton sensitivity --- measured through the real encoder at the real candidate width --- is
worth $+12.33$\pt{} of CaseHOLD accuracy at 3.25 bits on this model, at identical file size.
Pricing knapsack moves by the sensitivity \emph{difference} rather than by a score-times-params
proxy removes two hand-chosen constructs (the percentile rank and the $4^{-b}$ error curve) and
performs better than either.

\textbf{On the signals.} The gradient factor $\mathbb{E}[\delta^2]$ is the load-bearing one:
removing it flips the correlation with measured damage from $+0.521$ to $-0.338$. The activation
factor used alone anti-correlates ($\rho = -0.2059$), which means the published formula's second
term was subtracting information. Two per-channel second-moment vectors, sampled one step in
sixteen during a fine-tune that already happens, are enough.

\textbf{On allocation budgets.} Allocation only matters when the budget is tight. At 4.25 bits no
policy separates from uniform. At 3.25 bits the spread between policies is 40 points of accuracy
and even random allocation occasionally looks good. Any allocator claim made at a loose budget,
or without a random baseline, should be treated as unvalidated.

\textbf{On the runtime.} Packed execution is exact --- bit-identical accuracy to the simulated
path --- and the memory saving is real and predictable: 4.84\x{} peak allocated, with the
allocator's forecast within 0.03\pct{} of live. The kernel beats bf16 in isolation at every width
(1.09\x--2.56\x) and beats a non-templated quantized kernel by up to 2.36\x. It does not yet make
decode faster end to end, because the decode loop is host-bound at 22\pct{} GPU occupancy, and it
misses the 70\pct{} bandwidth gate at 2, 3 and 4 bits for a structural reason: at $K = 2048$ each
warp runs one main-loop iteration, so there is no intra-warp latency hiding, and the fp16
magic-number dequantization shortcut is unavailable to an fp32-accumulating kernel.

\textbf{The open question.} The premise that gave the method its name --- that the signals are
\emph{dynamic}, task-specific, so that a fine-tune on task $T$ yields an allocation specialised
to $T$ --- is the one thing not yet demonstrated. At 4.25 bits, two very different tasks produced
177 of 187 identical widths. Either the effect appears only under tight budgets, or the
allocation is largely a function of architecture and the ``dynamic'' framing needs revising. The
isolating experiment is small, well-defined, and next.

\vspace{1.5em}
\noindent{\color{dqblue}\rule{\textwidth}{0.8pt}}
\vspace{0.6em}

\noindent\small
\textbf{Provenance.} Every quantitative claim in this document traces to a measurement recorded in
\code{experiments/}\allowbreak\code{qwen35\_2b/}\allowbreak\code{RESULTS.md} in the DynQuant
repository, produced on the hardware and software stack named on the title page. Correlation figures are labelled with which ground-truth
key produced them (Section~\ref{sec:answerkey}); the two keys measure different quantities and
their numbers are not interchangeable. Statements about what is not established are in
Section~\ref{sec:notestablished} and are not hedges on the results above --- they are the boundary
of what was tested.

\end{document}
